\begin{Proposition}[{see \cite[Lemma~2 and Proposition~2]{ACGK}}]\label{properties_mutation_polygon}
Let the notation be the same as above.
\begin{itemize}\itemsep=0pt
\item[\rm (1)] If $Q\coloneqq\mutation_w(P,F)$, then we have $P=\mutation_{(-w)}(Q,F)$.
\item[\rm (2)] $P$ is a Fano polytope if and only if $\mutation_w(P,F)$ is a Fano polytope.
\end{itemize}
\end{Proposition}

Here, we recall that a convex lattice polytope $P\subset N_\RR$ with $\dim P=d$ is called \emph{Fano} if the origin is contained in the strict interior of $P$,
and the vertices $\calV(P)$ of $P$ are primitive lattice points of $N$.

Then, we consider the combinatorial mutation of a lattice polytope $P$ in terms of the polar dual $P^*$ of $P$ in $M$.
To do this, we must first discuss the polar dual $P^*$ for a polyhedron $P$.
We consider the family of polyhedra (not necessarily convex polytopes) which are of the following form:
\begin{displaymath}
\calP^d \coloneqq \bigg\{ \bigcap_{v \in S} H_{v,\geq -k_v} \cap \bigcap_{v' \in T}H_{v',\geq 0} \subset N_\RR \,|\, S,T \subset M, \; |S|, |T| < \infty, \; k_v \in \ZZ_{>0} \bigg\},
\end{displaymath}
where $H_{v,\geq k}=\{u \in N_\RR \,|\, \langle v,u \rangle \geq k\}$ for $v \in M$ and $k \in \RR$.
We note that any lattice polytope containing the origin of~$N_\RR$ belongs to $\calP^d$ but the ones not containing the origin do not belong to~$\calP^d$
since one of the supporting hyperplanes of such a polytope is of the form $H_{v, \geq k}$ for some $v \in M$ and some positive integer~$k$.
Similarly, we define~$\calQ^d$ by swapping the roles of~$N_\RR$ and~$M_\RR$.

For a given $P \in \calP^d$, we consider the \emph{polar dual} $P^* \subset M_\RR$ of $P$ defined as
\begin{displaymath}
P^* \coloneqq \{v \in M_\RR\,|\,\langle v,u\rangle\ge-1 \text{ for all } u\in P\}\subset M_\RR.
\end{displaymath}
Then, we have the following statements.

\begin{Proposition}\label{prop_dual_polytope}
Let the notation be the same as above. Then, for any $P \in \calP^d$ we have
\begin{itemize}\itemsep=0pt
\item[$(i)$] $P^* \in \calQ^d$,
\item[$(ii)$] $(P^*)^*=P$.
\end{itemize}
\end{Proposition}

\begin{proof}(i) Let $P \in \calP^d$. Since $P$ is a polyhedron, there exist a polytope~$Q$ and a polyhedral cone~$C$ such that
$P=Q+C$, where $+$ denotes the Minkowski sum (see \cite[Corollary~7.1b]{Sch}).
Let $Q=\operatorname{conv}\big(\big\{\frac{1}{k_1}u_1,\dots,\frac{1}{k_p}u_p\big\}\big)$ and let $C=\operatorname{cone}(\{u_1',\dots,u_q'\})$.
Note that we can choose~$u_i$,~$u_j'$ from $N$ and $k_i \in \ZZ_{>0}$ because of the form of~$P$.
In what follows, we will show that
\begin{displaymath}
P^*=\bigcap_{i=1}^p H_{u_i,\geq -k_i} \cap \bigcap_{j=1}^qH_{u_j',\geq 0}.
\end{displaymath}

First, we take $v \in P^*$. Then, we have $\langle v,u \rangle \geq -1$ for any $u \in P$.
Since $\frac{1}{k_i}u_i \in Q+{\bf 0} \subset P$, where ${\bf 0} \in C$ denotes the origin,
we see that $\langle v, \frac{1}{k_i}u_i \rangle \geq -1$, i.e., $\langle v, u_i \rangle \geq -k_i$ for each $i$.
If there is $j$ with $\langle v,u_j' \rangle<0$, then $\langle v, u'+ru_j' \rangle < -1$ for some $u' \in Q$ and some sufficiently large~$r$.
Moreover, we have $u'+ru_j' \in Q+C=P$. This contradicts $v \in P^*$; thus $\langle v,u_j' \rangle\geq 0$ for each~$j$.
Therefore, $v \in \bigcap_{i=1}^p H_{u_i,\geq -k_i} \cap \bigcap_{j=1}^qH_{u_j',\geq 0}$.

On the other hand, we take $v \in \bigcap_{i=1}^p H_{u_i,\geq -k_i} \cap \bigcap_{j=1}^qH_{u_j',\geq 0}$.
For any $u \in P$, as mentioned above, there exist $u' \in Q$ and $u'' \in C$ such that $u=u'+u''$.
Let $u'=\sum_{i=1}^p \frac{r_i}{k_i}u_i$, where $r_i \geq 0$ with $\sum_{i=1}^pr_i=1$, and let $u''=\sum_{j=1}^q s_ju_j'$, where $s_j \geq 0$.
By using these expressions together with the inequalities $\langle v,u_i \rangle \geq -k_i$ for each $i$ and $\langle v,u_j' \rangle \geq 0$ for each $j$, we see that
\begin{align*}
\langle v,u \rangle = \langle v,u' \rangle+\langle v,u'' \rangle =
\sum_{i=1}^p \frac{r_i}{k_i}\langle v,u_i \rangle+\sum_{j=1}^q s_j\langle v,u_j' \rangle \geq -\sum_{i=1}^pr_i=-1,
\end{align*}
and thus we have $v \in P^*$.

(ii) For any $u \in P$, we have $\langle v,u\rangle\geq -1$ for any $v \in P^*$, which means that $P \subset (P^*)^*$.
For the other inclusion, we take $u \in N_\RR \setminus P$.
Let $P=\bigcap_{v \in S} H_{v,\geq -k_v} \cap \bigcap_{v' \in T}H_{v',\geq 0}$.
Then either $\langle v,u \rangle <-k_v$ for some $v \in S$ or $\langle v',u \rangle < 0$ for some $v' \in T$ holds.
In the former case, since $\langle v,u' \rangle \geq -k_v$ for any $u' \in P$, we have $\frac{1}{k_v}v \in P^*$.
This means that there is $v''\coloneqq\frac{1}{k_v}v \in P^*$ such that $\langle v'',u \rangle<-1$, and hence $u \not\in (P^*)^*$.
Similarly, in the latter case, since $\langle rv',u' \rangle \geq 0 \geq -1$ for any $u' \in P$ and $r \geq 0$, we have $rv' \in P^*$.
This implies that $u \not\in (P^*)^*$ for sufficiently large $r$, and hence $u \not\in (P^*)^*$. Therefore, we obtain $(P^*)^* \subset P$, as required.
\end{proof}

We next define a map
\[
\varphi_{w,F}\colon \ M_\RR\rightarrow M_\RR \qquad \text{as} \quad \varphi_{w,F}(v)\coloneqq v-v_{\min}w,
\]
where \mbox{$v_{\min}\coloneqq {\min}\{\langle v, u\rangle\,|\, u\in F\}$}.
In particular, when $d=2$ (see Remark~\ref{rem_def_mutation}), this map can be described as
\begin{equation}\label{mutation_M}
\varphi_{w,F}(v)=
\begin{cases}
v & \text{if $\langle v,u_E\rangle\ge0$}, \\
v-\langle v,u_E\rangle w&\text{if $\langle v,u_E\rangle<0$},
\end{cases}
\end{equation}
with $F=\operatorname{conv}\{{\bf 0},u_E\}$.
The next proposition is crucial to prove our main result Theorem~\ref{mutation=deformation}.


\begin{Proposition}
\label{prop_mutation_dual_polytope}
For any $P \in \calP^d$, we have
\begin{displaymath}\varphi_{w,F}(P^*)=\mut_w(P,F)^*.
\end{displaymath}
\end{Proposition}

\begin{proof}
Although this equality essentially follows from \cite[Proposition~4]{ACGK} and the discussions in \cite[p.~12]{ACGK}, we give a precise proof for completeness.

Let $\varphi=\varphi_{w,F}$ and $Q=\mut_w(P,F)$.
To show $\varphi(P^*)\subset Q^*$, we take $v \in P^*$ arbitrarily and consider $\varphi(v)=v - v_{\min}w \in \varphi(P^*)$.
We will show that $\langle v - v_{\min}w,u \rangle \geq -1$ for any $u \in Q$.
It suffices to show this for each vertex $u\in\calV(Q)$.
\begin{itemize}\itemsep=0pt
\item Let $u\in\calV(Q)$ with $\langle w,u \rangle \geq 0$. Then we can write $u=u_P+\langle w,u_P \rangle u_F$ for some $u_P \in \calV(P)$ and $u_F \in \calV(F)$.
In particular, we have
\begin{displaymath}
\langle v-v_{\min}w, u\rangle = \langle v,u_P \rangle + \langle w,u_P \rangle(\langle v,u_F \rangle-v_{\min}) \geq \langle v,u_P \rangle \geq -1.
\end{displaymath}

\item Let $u\in\calV(Q)$ with $\langle w,u \rangle < 0$. For any $u_F \in \calV(F)$, we have $u-\langle w,u \rangle u_F \in P$.
Hence, $\langle v,u-\langle w,u \rangle u_F \rangle \geq -1$. In particular, $\langle v,u \rangle \geq -1 + v_{\min} \langle w,u \rangle$. Thus, we see that
\begin{displaymath}
\langle v-v_{\min}w, u\rangle =\langle v, u\rangle - v_{\min} \langle w, u\rangle \geq -1.
\end{displaymath}
\end{itemize}

To show $Q^*\subset\varphi(P^*)$, we will show that for any $v \in Q^*$ there is $v' \in P^*$ such that $v=\varphi(v')$.
Let $\Delta_F$ be the normal fan of $F$ in $M_\RR$ and let $\sigma \in \Delta_F$ be a maximal cone in $\Delta_F$.
The discussions in \cite[p.~12]{ACGK} say that there exists $M_\sigma \in \GL_d(\ZZ)$ such that the map $\varphi$ is equal to $M_\sigma$,
i.e., $\varphi(v)=vM_\sigma$. Thus, we conclude that $\varphi(v')=v$ for $v'=vM_\sigma^{-1}$.
\end{proof}
